\documentclass[a4j,dvipdfmx,10pt]{ujreport}
\usepackage{amsmath,amssymb,amsthm}
\usepackage{amsmath,amssymb}

\usepackage{mathrsfs}
%筆記体や花文字を出力するために必要

%図表関係
\usepackage{tabularx}
\usepackage{longtable}
%ページをまたぐ表を出力するために必要
\usepackage{multirow}
%表のセルを結合するために必要
\usepackage{here}
%図表を指定した位置に出力するのに必要
\renewcommand{\tablename}{Table}
%表下に表示される「表：（番号）」が「table:（番号）」になる
\usepackage{colortbl}
%表の背景に色を塗れるようになる
%↓はそのために定義した色
\definecolor{pink}{cmyk}{0.093, 0.551, 0.289, 0}

\usepackage[dvipdfmx]{graphicx}

%画像を回り込み文章を書くために必要
\usepackage{wrapfig}

\usepackage{ascmac}
%囲み付き文章を出力するために必要

\usepackage{comment}
%複数行をコメントアウトするために必要

\usepackage{url}
%urlを出力するために必要

\usepackage{makeidx}
%索引を作るために必要

\usepackage{subfiles}
%ファイルを小分けにして編集するために必要

\usepackage{pifont}
%指差しマーク出力のために

%becouse-in-proof環境の背景に色を付けるために必要なもの
\usepackage{tcolorbox}
\tcbuselibrary{breakable, skins, theorems}
\usepackage{enumitem}
%itemize環境の細かい余白制御に使う
%自前のなぜならば環境を定義するときに使っている

%TikZ使うためのもの。documentclassを\documentclass[a4j,dvipdfmx,10pt]{jbook}にしておくこと。
\usepackage{tikz}
\usetikzlibrary{intersections, calc, arrows}
\usepackage{tikz-qtree}

%ノート内の参照やURLにハイパーリンクが埋め込めるようになる
\usepackage[dvipdfmx,bookmarkstype=toc,colorlinks=true,urlcolor=black,%
      linkcolor=blue,citecolor=blue,linktocpage=true,bookmarks=true]{hyperref}
      \AtBeginDvi{\special{pdf:tounicode 90ms-RKSJ-UCS2}}


%見栄え関係
%余白の設定。各cm分余白が削られる
\usepackage{geometry}
\geometry{top=2cm, bottom=2cm, left=1cm, right=1cm, includefoot}

%見出しなどの設定
\newtheoremstyle{mystyle}%   % スタイル名
    {5pt}%                      % 上部スペース
    {12pt}%                      % 下部スペース
    {\normalfont\setlength{\parindent}{2em}}%           % 本文フォント
    {}%                      % インデント量
    {\bf}%                   % 見出しフォント
    {.}%                      % 見出し後の句読点, '.'
    {\newline}%                     % 見出し後のスペース, ' ' or \newline
    {\underline{\thmname{#1}{\;}\thmnumber{#2}}\thmnote{（#3）}}


\theoremstyle{mystyle}

\newtheorem{my-theorem}{Theorem}[section]
\newtheorem{my-lemma}[my-theorem]{Lemma}
\newtheorem{my-proposition}[my-theorem]{Proposition}
\newtheorem{my-corollary}[my-theorem]{Corollary}
\newtheorem{my-definition}[my-theorem]{Definition}
\newtheorem{my-example}[my-theorem]{Example}
\newtheorem{my-remark}[my-theorem]{Remark}
\newtheorem{my-exercise}[my-theorem]{Exercise}

%ここから全ての定理環境を上書きする.
%これで主張終わりに黒四角が出力され,\qedhereも使えるようになる.
\newcommand{\thmsymbol}{$\blacksquare$}

\newenvironment{theorem}[1][]{
  \begin{my-theorem}[#1]%
  \renewcommand{\qedsymbol}{\thmsymbol}\pushQED{\qed}}%
  {\popQED\end{my-theorem}
}
\newenvironment{lemma}[1][]{
  \begin{my-lemma}[#1]%
  \renewcommand{\qedsymbol}{\thmsymbol}\pushQED{\qed}}%
  {\popQED\end{my-lemma}
}
\newenvironment{proposition}[1][]{
  \begin{my-proposition}[#1]%
  \renewcommand{\qedsymbol}{\thmsymbol}\pushQED{\qed}}%
  {\popQED\end{my-proposition}
}
\newenvironment{corollary}[1][]{
  \begin{my-corollary}[#1]%
  \renewcommand{\qedsymbol}{\thmsymbol}\pushQED{\qed}}%
  {\popQED\end{my-corollary}
}
\newenvironment{definition}[1][]{
  \begin{my-definition}[#1]%
  \renewcommand{\qedsymbol}{\thmsymbol}\pushQED{\qed}}%
  {\popQED\end{my-definition}
}
\newenvironment{example}[1][]{
  \begin{my-example}[#1]%
  \renewcommand{\qedsymbol}{\thmsymbol}\pushQED{\qed}}%
  {\popQED\end{my-example}
}
\newenvironment{remark}[1][]{
  \begin{my-remark}[#1]%
  \renewcommand{\qedsymbol}{\thmsymbol}\pushQED{\qed}}%
  {\popQED\end{my-remark}
}
\newenvironment{exercise}[1][]{
  \begin{my-exercise}[#1]%
  \renewcommand{\qedsymbol}{\thmsymbol}\pushQED{\qed}}%
  {\popQED\end{my-exercise}
}

%ここから証明環境の構築
\renewcommand{\proofname}{Proof}

\makeatletter % use at mark
\renewenvironment{proof}[1][\proofname]{\par
  \pushQED{\qed}%
  \normalfont \topsep6\p@\@plus6\p@\relax
  \trivlist
  \item[\hskip\labelsep
        \itshape
    {\bf\underline{#1}\@addpunct{\ }}]\ignorespaces
    % {\bf\underline{#1}\@addpunct{.}}]\ignorespaces % ピリオドあり
}{%
  \popQED\endtrivlist\@endpefalse
}
\makeatother % end at mark

%証明と同じ処理をする別の解答環境
\newenvironment{answer}{\begin{proof}[Answer]}{\end{proof}}

%マクロ

%%各種参照
\newcommand{\PartRef}[1]{第\ref{#1}部（\pageref{#1}ページ）}
\newcommand{\ChapRef}[1]{第\ref{#1}章（\pageref{#1}ページ）}
\newcommand{\SecRef}[1]{\ref{#1}節（\pageref{#1}ページ）}
\newcommand{\DefRef}[1]{Definition \ref{#1}（\pageref{#1}ページ）}
\newcommand{\ExRef}[1]{Example \ref{#1}（\pageref{#1}ページ）}
\newcommand{\NoteRef}[1]{Notation \ref{#1}（\pageref{#1}ページ）}
\newcommand{\LemRef}[1]{Lemma \ref{#1}（\pageref{#1}ページ）}
\newcommand{\ThRef}[1]{Theorem \ref{#1}（\pageref{#1}ページ）}
\newcommand{\FactRef}[1]{Fact \ref{#1}（\pageref{#1}ページ）}
\newcommand{\PropRef}[1]{Proposition \ref{#1}（\pageref{#1}ページ）}
\newcommand{\PuzzleRef}[1]{Puzzle \ref{#1}（\pageref{#1}ページ）}
\newcommand{\TableRef}[1]{表 \ref{#1}}

%%証明中の部分証明環境
\definecolor{BIP-back}{gray}{0.85}
%sub-block の入れ子深さを数える
\newcount\subblocklevel
\subblocklevel=0

\newcommand{\subblockcolback}{%
  \ifodd\subblocklevel BIP-back\else white\fi
}

\makeatletter
% sub-block 内専用の脚注：箱の外へ回す
\newcommand{\sbfootnote}[1]{%
  \footnotemark
  \expandafter\g@addto@macro\csname sb@footnotes@\the\subblocklevel\endcsname{%
    \footnotetext{#1}%
  }%
}
\makeatother

\newenvironment{sub-block}[1]
  {%
    \global\advance\subblocklevel by 1\relax
    \expandafter\gdef\csname sb@footnotes@\the\subblocklevel\endcsname{}%
    \vspace{-5pt}
    \begin{tcolorbox}[
      colback = \subblockcolback,
      colframe = white,
      boxrule = 5pt,
      fonttitle = \bfseries,
      breakable = true]
    \begin{itemize} \item[#1]
  }
  {%
    \end{itemize}
    \end{tcolorbox}
    \vspace{-5pt}
    % ←ここで「箱の外」に脚注本文を出す
    \csname sb@footnotes@\the\subblocklevel\endcsname
    \global\advance\subblocklevel by -1\relax
  }


%%定義する両矢印
\newcommand{\defarr}{\ \mbox{$\stackrel{\text{def}}{\Longleftrightarrow}$}}
\newcommand{\tempdefarr}{:\Leftrightarrow}
%%定義する等号
\newcommand{\defeq}{\ \mbox{$\stackrel{\text{def}}{=}$}}
\newcommand{\tempdefeq}{:=}

%%同値性証明を書く時の略記
\newcommand{\ImpDir}[2]{
  \hspace{-0.5cm}（\;(#1)$\Rightarrow$(#2)\;）
}

\newcommand{\dash}[1]{#1 ^{\prime}}
%集合関連
%%集合の内包的記法
\newcommand{\SetBar}[2]{
  \{\,#1\,\mid\,#2\,\}
}
\newcommand{\SetColon}[2]{
  \{\,#1\,:\,#2\,\}
}
%%順序対
\newcommand{\op}[2]{\langle #1,#2 \rangle}
%写像関連
\newcommand{\map}[3]{#1\colon #2 \to #3}
\newcommand{\InjectionMap}[3]{
  #1\colon #2 \xrightarrow{\text{inj}} #3
}
\newcommand{\SurjectionMap}[3]{
  #1\colon #2 \xrightarrow{\text{sur}} #3
}
\newcommand{\BijectionMap}[3]{
  #1\colon #2 \xrightarrow{\text{bij}} #3
}
\DeclareMathOperator{\dom}{dom}
\DeclareMathOperator{\ran}{ran}
%%恒等関数
\newcommand{\idmap}[1]{\mathrm{id}_{#1}}
%%濃度
\DeclareMathOperator{\card}{card}

\title{集合論濃度編ゼミノート}
\author{静間}
\date{\today}

\begin{document}

\maketitle

レッスンのカリキュラム内の定義や記法と異なる点のメモ
\begin{enumerate}
  \item 
    同値の意味で「$\Leftrightarrow$」を使っている.
    レッスンでは単なる文章や論理式の言い換えに使っていた.
  \item レッスンでいう$\InjectionMap{f}{X}{Y}$のような単射・全射を表す記法がない？
\end{enumerate}

\newpage

\setcounter{chapter}{5}
\chapter{有限と無限}
\label{chapter:有限と無限}
%テキストの章番号と一致するように設定

\setcounter{section}{0}
\section{集合の対等関係}
\label{section:集合の対等関係}
%テキストの節番号と一致するように設定

\begin{definition}[集合の対等関係, p66の定義6.1]
  \label{definition:集合の対等関係}
  2つの集合$A,B$に対して$A$から$B$への全単射写像が存在するとき,
  $A$と$B$は\textgt{対等}であるといい,
  $A\sim B$で表す.
  つまり$A\sim B \defarr \exists f( \; \BijectionMap{f}{X}{Y} \; )$.
\end{definition}

\begin{example}[p66の例6.2]
  $A\tempdefeq \{ 0,1,2 \},B\tempdefeq \{ a,b,c \}$とおくとき,
  $A\sim B$である.
  なぜならば$f$を$\map{f}{A}{B}$で$f(0)=a,f(1)=b,f(2)=c$なるものとして定めると,
  $f$は全単射であり,
  そんな$f$の存在から$A\sim B$であることが分かる.
\end{example}

\begin{boxnote}
  記号「\ding{43}」のテキストにおける用法を調べておく.
\end{boxnote}

\begin{proposition}
  1つ上の例に出てきた$f$は全単射である.
\end{proposition}

\begin{proof}
  葉山さんにパス.
\end{proof}

\begin{proposition}[p66]
  有限集合$A$の要素の数を$|A|$で表すとする. \\
  有限集合$X,Y$について以下は同値.
  \begin{itemize}
    \item[(a)] $|X|=|Y|$.
    \item[(b)] $X\sim Y$. \qedhere
  \end{itemize}
\end{proposition}

\begin{proof}
  \mbox{}
  \begin{itemize}
    \item[] \ImpDir{a}{b} \\
      $X$は有限集合なので$n$を$|X|=n$な$n\in \mathbb{N}$とする.
      よって$X$を$\{ x_1,\dots,x_n \}$と表す.
      同様に$Y$を$\{ y_1,\dots,y_n \}$と表す.
      $\map{f}{X}{Y}$な$f$を$i\in \{ 1,\dots,n\}$について$f(x_i)=y_i$で定める.
      順序対の集合として$f$を表すならば$\SetBar{ \op{x_i}{y_i} }{ i\in \{ 1,\dots,n\} }$となる.
      ここで$f$は全単射である.
      \begin{sub-block}{{\boldmath $\because$}\,}
        \begin{itemize}
          \item[] \hspace{-0.5cm}（単射性） \\
            任意に$x\neq \dash{x}$な$x,\dash{x}\in X$をとる.
            するとある$i,j\in \{ 1,\dots,n\}$があって,
            $i\neq j\land x=x_i\land \dash{x}=x_j$となっているので,
            そんな$i,j$をそれぞれとって固定しておく.
            すると$f(x)=f(x_i)=y_i$と$f(\dash{x})=f(x_j)=y_j$となるが,
            $i\neq j$より$y_i\neq y_j$,
            つまり$f(x)\neq f(\dash{x})$である.

          \item[] \hspace{-0.5cm}（全射性） \\ 
            任意に$y\in Y$をとる.
            するとある$i\in \{ 1,\dots,n\}$があって,
            $y=y_i$となっているので,
            そんな$i$をとって固定しておく.
            ここで$x\tempdefeq x_i$とおくと$x\in X$かつ$f(x)=f(x_i)=y_i=y$より$f(x)=y$である.
            そんな$x$の存在より$\exists x\in X( f(x)=y )$は真である.
        \end{itemize}
      \end{sub-block}
      そんな$f$の存在より$X\sim Y$である.
      
    \item[] \ImpDir{b}{a} \\
      $X\sim Y$より存在する$\BijectionMap{f}{X}{Y}$な$f$を1つとり固定する.
      $|X|=|f[X]|$である.
      \begin{sub-block}{{\boldmath $\because$}\,}
        $X$が有限集合なので,
        $f[X]$も有限である.
        \begin{sub-block}{{\boldmath $\because$}\,}
          $X$が有限なので$n\in \mathbb{N}$を用いて$|X|=n$とおくと,
          $X$は$\{x_1,\dots , x_n\}$と表せて,
          $f[X]$は$\{ f(x_1),\dots,f(x_n)\}$と表せることからその要素数は明らかに$n$個以下であり,
          ゆえにその要素数は有限である.
        \end{sub-block}

        そしていつでも$|f[X]|\leq |X|$が成立する.
        \begin{sub-block}{{\boldmath $\because$}\,}
          $|f[X]|> |X|$だったとする.
          $n,n_f\in \mathbb{N}$を用いて
          $|f[X]|=n_f,|X|=n$とおくと,
          $X$は$\{x_1,\dots , x_n\}$と表せて,
          $f[X]$は$\{ f(x_1),\dots,f(x_n)\}$と表せる.
          そして$f[X]$の要素の個数は$n$以下であるが,
          $|f[X]|=n_f$よりその要素数が$n$より大きい$n_f$であることは矛盾である.
        \end{sub-block}

        ここで背理法を用いるため$|f[X]|\neq |X|$だったとすると,
        $|f[X]|< |X|$である.
        $n,n_f\in \mathbb{N}$を用いて
        $|f[X]|=n_f,|X|=n$とおくと,
        $X$は$\{x_1,\dots , x_n\}$と表せて,
        $f[X]$は$\{ f(x_1),\dots,f(x_n)\}$と表せる.
        ここで$|f[X]|=n_f<n$より$f(x_1),\dots,f(x_n)$たちにおいて,
        ある$i,j(\i\neq j)$があって$f(x_i)=f(x_j)$となってなくてならない
        \sbfootnote{
          これは\textgt{鳩ノ巣原理}と呼ばれている.
          ざっくりと説明すると$n<k$な自然数$n,k$があって$k$個のものを$n$個のグループに分けるとき,
          どこかのグループには2つ以上のものが所属することになるということである.
          言い換えると$k$個のものたちを1個ずつ所属するグループに分けるならば,
          そのグループは$k$と同じ数だけ必要であるとことである.
          この原理は色々な表現方法がある.
          この数学教室のお客様であれば\cite{Basic001}のp143も参考にしてください.
        }.

        そのような$x_i,x_j$の存在は$f$が単射であったことに矛盾する.
        よって$|X|=|f[X]|$.
      \end{sub-block}

      また$f$が全射なので$f[X]=Y$なので$|f[X]|=|Y|$.
      よって$|X|=|Y|$. \qedhere
  \end{itemize}
\end{proof}

\begin{example}[p67の例6.3]
  \label{example:自然数全体と偶数全体は対等}
  偶数の集合$E:=\SetBar{ 2n }{ n\in \mathbb{N} }$に対して$\mathbb{N}\sim E$が成立する.
  なぜならば$\map{f}{\mathbb{N}}{E}$な$f$を$f(n)=2n$と定めれば$f$は$\mathbb{N}$から$E$への全単射であり,
  そんな$f$の存在から$\mathbb{N}\sim E$である.
\end{example}

\begin{proposition}
  Example \ref{example:自然数全体と偶数全体は対等}における$f$は$\mathbb{N}$から$E$への全単射である.
\end{proposition}

\begin{proof}
  葉山さんにパス.
\end{proof}

\begin{example}[p67の例6.4]
  \label{example:自然数全体と整数全体は対等}
  $\mathbb{N}\sim \mathbb{Z}$が成立する.
  なぜならば$\map{f}{\mathbb{N}}{\mathbb{Z}}$な$f$を
  \begin{equation*}
    f(n)=
    \begin{cases}
      n/2 & \text{（}n\text{が偶数のとき）},\\
      (1-n)/2  & \text{o.w.}
    \end{cases}
  \end{equation*}
  と定めれば$f$は$\mathbb{N}$から$\mathbb{Z}$への全単射であり,
  そんな$f$の存在から$\mathbb{N}\sim \mathbb{Z}$である.
\end{example}

\begin{proposition}
  \label{proposition:自然数全体と整数全体は対等}
  Example \ref{example:自然数全体と整数全体は対等}における$f$は$\mathbb{N}$から$\mathbb{Z}$への全単射である.
\end{proposition}

\begin{proof}
  \mbox{}
  \begin{itemize}
    \item[] \hspace{-0.5cm}（単射性） \\
      任意に$n_1,n_2\in \mathbb{N}$をとり$n_1\neq n_2$とする.
      $n_1,n_2$の偶奇が同じかどうかで場合分けする.
      \begin{itemize}
        \item[] \hspace{-0.5cm}（$n_1,n_2$の偶奇が同じとき） \\
          $n_1,n_2$がともに偶数だったとする.
          つまり$m_1,m_2\in \mathbb{N}$があって$n_1=2m_1,n_2=2m_2$となっている.
          ここで$n_1\neq n_2$より$m_1\neq m_2$である.
          そして$f(n_1)=f(2m_1)=2m_1/2=m_1$と$f(n_2)=f(2m_2)=2m_2/2=m_2$より$f(n_1)\neq f(n_2)$である. \\
          $n_1,n_2$がともに奇数であったときも同様に示せるので割愛する.

        \item[] \hspace{-0.5cm}（$n_1,n_2$の偶奇が異なるとき） \\
          $n_1$が偶数,
          $n_2$が奇数であったとしても一般性を失わない.
          すると$f(n_1)\geq 0,f(n_2)<0$より$f(n_1)\neq f(n_2)$である.

      \end{itemize}

    \item[] \hspace{-0.5cm}（全射性） \\
      任意に$y\in \mathbb{Z}$をとる.
      $y\geq 0$のとき$x\tempdefeq 2y$とおくと,
      $x\in \mathbb{N}$であり偶数である.
      そして$f(x)=f(2y)=y$となって,
      そんな$x$の存在より$\exists x\in \mathbb{N}( f(x)=y )$は真である. \\
      $y<0$のときは$x\tempdefeq 1-2y$とおけば同様に示せる.
    
  \end{itemize}
\end{proof}

\begin{example}[p67の例6.5]
  \label{example:任意の閉区間同士は対等}
  任意の2つの閉区間$I\tempdefeq [a,b],\dash{I}=[c,d](a<b,c<d)$について$I\sim \dash{I}$が成立する.
  なぜならば$\map{f}{I}{\dash{I}}$な$f$を
  \begin{equation*}
    f(x)=\frac{c-d}{a-b}x+\frac{ad-bc}{a-b}
  \end{equation*}
  と定めれば$f$は$I$から$\dash{I}$への全単射であり,
  そんな$f$の存在から$I\sim \dash{I}$である.
\end{example}

\begin{proposition}
  Example \ref{example:任意の閉区間同士は対等}における$f$は$I$から$\dash{I}$への全単射である.
\end{proposition}

\begin{proof}
  葉山さんにパス.
\end{proof}

\begin{example}[p67の例6.6]
  \label{example:開区間と実数全体は対等}
  開区間$J\tempdefeq (-1,1)$に対して$\mathbb{R}\sim J$が成立する.
  なぜならば$\map{f}{\mathbb{R}}{J}$な$f$を
  \begin{equation*}
    f(x)=\frac{x}{1-|x|}
  \end{equation*}
  と定めれば$f$は$\mathbb{R}$から$J$への全単射であり,
  そんな$f$の存在から$\mathbb{R}\sim J$である.
\end{example}

\begin{proposition}
  Example \ref{example:開区間と実数全体は対等}における$f$は$\mathbb{R}$から$J$への全単射である.
\end{proposition}

\begin{proof}
  葉山さんにパス.
\end{proof}

\begin{exercise}[p67の問1]
  $2$を加えると$5$の倍数になる整数全体の集合$A$に対し$\mathbb{N}\sim A$が成り立つことを示せ.
\end{exercise}

\begin{answer}
  $A$は$\SetBar{ 5x-2 }{ x\in \mathbb{Zb} }$と表せる（外延的記法ならば$\{ \dots,-12,-7,-2,3,8,\dots \}$である）.
  ここで$\map{f}{\mathbb{N}}{A}$な$f$を$n\in \mathbb{N}$に対して
  \begin{equation*}
    f(n)=
    \begin{cases}
      5\cdot n/2 -2 & \text{（}n\text{が偶数のとき）},\\
      5\cdot (1-n)/2 -2  & \text{o.w.}
    \end{cases}
  \end{equation*}
  で定めると,
  $f$は$\mathbb{N}$から$A$への全単射である（その証明は葉山さんにパス）.
  そんな$f$の存在より$\mathbb{N}\sim A$.
\end{answer}

\noindent
対等であることを示すとき,
ちょうどよい1つの全単射写像を思いつかなければProposition \ref{proposition:対等関係は同値関係}にある性質を利用して,
いくつか全単射を経由して証明する方法もある.
上の回答をあえてProposition \ref{proposition:対等関係は同値関係}を使うやり方で示してみる.

\begin{answer}
  $\map{g}{\mathbb{Z}}{A}$な$g$を$x\in \mathbb{Z}$に対して$g(x)=5x-2$で定めると,
  （入力を定数倍して定数を引くだけというレベルで）Proposition \ref{proposition:自然数全体と整数全体は対等}と同様にして$g$は全単射であることがわかる.
  よってそんな$g$の存在から$\mathbb{Z}\sim A$である.
  Example \ref{example:自然数全体と整数全体は対等}より$\mathbb{N}\sim \mathbb{Z}$であるから$\mathbb{Z}\sim A$と合わせて,
  Proposition \ref{proposition:対等関係は同値関係}の推移律より$\mathbb{N}\sim A$.
\end{answer}

\begin{exercise}[p68の問2]
  任意の半開区間$H\tempdefeq [a,b),\dash{H}=[c,d)(a<b,c<d)$は対等であることを示せ.
\end{exercise}

\begin{answer}
  葉山さんにパス.
\end{answer}

\begin{exercise}[p68の問3]
  区間$J\tempdefeq (0,+\infty)$に対し,
  $\mathbb{R}\sim J$が成り立つことを示せ.
\end{exercise}

\begin{answer}
  葉山さんにパス.
\end{answer}

\begin{theorem}[Cantor, p68の定理6.7]
  \label{theorem:自然数全体と実数全体は対等でない}
  $\mathbb{N} \nsim \mathbb{R}$.
\end{theorem}

\begin{proof}
  葉山さんにパス.
\end{proof}

\begin{exercise}[p69の問4]
  Theorem \ref{theorem:自然数全体と実数全体は対等でない}の証明における$b_n$の定め方について,
  なぜ$b_n=9$とすると不都合であるか考えよ.
\end{exercise}

\begin{answer}
  葉山さんにパス.
\end{answer}

\begin{proposition}[p69の命題6.8]
  \label{proposition:対等関係は同値関係}
  任意の集合$A,B,C$に対して,
  次が成り立つ.
  \begin{enumerate}
    \item $A\sim A$.（反射律）
    \item $A\sim B$ならば$B\sim A$.（対称律）
    \item $A\sim B \land B\sim C$ならば$A\sim C$.（推移律） \qedhere
  \end{enumerate}
\end{proposition}

\begin{proof}
  \mbox{}
  \begin{enumerate}
    \item 
      $A$上の恒等写像を$\idmap{A}$と表せば,
      $\map{\idmap{A}}{A}{A}$であり,
      $\idmap{A}$は全単射である.
      そんな写像の存在から$A\sim A$.

    \item 
      $A\sim B$より存在する$\BijectionMap{f}{A}{B}$な$f$を1つとり固定する.
      $\BijectionMap{f}{A}{B}$より$f$の逆写像$f^{-1}$が存在し,
      $\BijectionMap{f^{-1}}{B}{A}$である.
      そんな写像の存在より$B\sim A$.

    \item 
      $A\sim B$と$B\sim C$より存在する$\BijectionMap{f}{A}{B},\BijectionMap{g}{B}{C}$をそれぞれ1つずつとり固定する.
      その合成$g\circ f$は$\BijectionMap{g\circ f}{A}{C}$であり,
      そんな写像の存在より$A\sim C$. \qedhere
  \end{enumerate}
\end{proof}

\renewcommand{\bibname}{参考文献}
%設定しないと「関連図書」になる

\bibliographystyle{jplain}
\bibliography{souji.bib}

%\printindex
%索引を出力

\end{document}
